% Four Conics Theorem
% https://mathworld.wolfram.com/FourConicsTheorem.html
% If two intersections of each pair of three conics 
% S_1, S_2, and S_3 lie on a conic S_0, then the 
% lines joining the other two intersections of each 
% pair are concurrent (Evelyn et al. 1974, pp. 23 and 25). 

\documentclass{standalone}
\usepackage{tikzplot}

\begin{document}

\begin{tikzpicture}
  \tikzmath{
    \A{1,1} = -2; \A{2,1} = 3; \A{3,1} = 1;
    \B{1,1} = 1; \B{2,1} = -3; \B{3,1} = 1; 
    \C{1,1} = 1; \C{2,1} = 3; \C{3,1} = 1; 
    \D{1,1} = -2; \D{2,1} = -3; \D{3,1} = 1; 
    \E{1,1} = 3; \E{2,1} = -1; \E{3,1} = 1; 
    % \F{1,1} = -3; \F{2,1} = 1; \F{3,1} = 1; 
  }
  \tikzset{
    % Create a conic through five given points (A to E)
    conics/define/five points={\A,\B,\C,\D,\E,\Co},
    % Create the sixth point F
    conics/midpoint={\C,\D,\M},
    conics/intersect cl={\Co,\E,\M,\P,\Q},
    conics/exclude={\P,\Q,\E,\F},
    % Create three conics
    conics/define/pencil={\A,\B,\C,\D,-3,\Ca},
    conics/define/pencil={\A,\B,\E,\F,-3,\Cb},
    conics/define/pencil={\C,\D,\E,\F,-10,\Cc},
    % line AB and midpoint
    conics/cross={\A,\B,\lab},
    conics/midpoint={\A,\B,\Mab},
    % line CD and midpoint
    conics/cross={\C,\D,\lcd},
    conics/midpoint={\C,\D,\Mcd},
    % line EF and midpoint
    conics/cross={\E,\F,\lef},
    conics/midpoint={\E,\F,\Mef},
    % other line of intersections
    conics/linear factor={\Ca,\Cb,\lab,\Mab,\l},
    conics/linear factor={\Ca,\Cc,\lcd,\Mcd,\m},
    conics/linear factor={\Cb,\Cc,\lef,\Mef,\n},
  }

  \clip (-6,-6) rectangle (6,6);
  \axes[grid] (-5:5,-5:5);
  \conic[dashed,draw=violet] (\Co);
  \conic[thick,draw=teal] (\Ca);
  \conic[thick,draw=magenta] (\Cb);
  \conic[thick,draw=purple] (\Cc);
  \segment[thick,draw=navy,clip=(-5:3,-4:4)] (\l);
  \segment[thick,draw=teal,clip=(-5:3,-5:4)] (\m);
  \segment[thick,draw=red,clip=(-4:2,-5:4)] (\n);

  \coordinate (A) at (\A{1,1},\A{2,1});
  \coordinate (B) at (\B{1,1},\B{2,1});
  \coordinate (C) at (\C{1,1},\C{2,1});
  \coordinate (D) at (\D{1,1},\D{2,1});
  \coordinate (E) at (\E{1,1},\E{2,1});
  \coordinate (F) at (\F{1,1}/\F{3,1},\F{2,1}/\F{3,1});

  \foreach \p/\placement in {A/above left, B/below right,
  C/above,D/above,E/above,F/above}{
    \fill[red] (\p) circle (1.5pt);
    \draw (\p) node[\placement] {$\p$};
  }
\end{tikzpicture}

\end{document}